\section{Архитектура на системата}

След формулирането на изискванията естествената следваща стъпка е анализ на архитектурната организация на приложението. Архитектурата е гръбнакът на всяка нетривиална система, защото определя как отделните отговорности са разпределени, как се движат данните между модулите и до каква степен проектът може да се поддържа и развива във времето. В настоящата разработка архитектурният избор е насочен към ясно разделение между уеб слой, бизнес логика, достъп до данни, конфигурация и визуализация. Тази глава описва именно това разпределение и го свързва с реалната структура на сорс кода.

\subsection{Обща архитектурна постановка}

Системата е реализирана като монолитно уеб приложение върху Spring Boot. Терминът „монолитно“ тук не следва да се възприема като негативен. В контекста на настоящия проект той означава, че всички основни функционални компоненти -- уеб маршрутизация, бизнес логика, persistence слой, HTML шаблони и конфигурация -- се намират в един проект и се стартират в един процес. Това решение е подходящо поради няколко причини.

Първо, домейнът на системата е относително компактен. Няма независими бизнес домейни, които да изискват отделни услуги или самостоятелни deployment единици. Второ, учебната цел на проекта е да демонстрира завършена, но прозрачна архитектура, а не да усложнява излишно инфраструктурния модел. Трето, Spring Boot предоставя достатъчно добра организационна рамка, така че монолитното приложение да остане вътрешно добре структурирано.

Архитектурната схема може да бъде представена като съвкупност от следните слоеве:

\begin{itemize}[leftmargin=*]
    \item слой за HTTP комуникация и управление на заявки;
    \item слой за бизнес операции и домейн логика;
    \item слой за достъп до данните и persistence;
    \item слой за представяне чрез HTML шаблони;
    \item конфигурационен слой за bean компоненти, локализация и инфраструктурни настройки.
\end{itemize}

\begin{figure}[H]
    \centering
    \begin{tikzpicture}
        \node[appblock] (client) {Клиентски браузър};
        \node[appblock, below=1.3cm of client] (web) {Web слой\\контролери и маршрути};
        \node[appblock, below=1.3cm of web] (service) {Service слой\\бизнес операции};
        \node[appblock, below=1.3cm of service] (repo) {Repository слой\\JPA интерфейси};
        \node[appblock, below=1.3cm of repo] (db) {MySQL};
        \node[appblock, right=4.2cm of web] (view) {Thymeleaf шаблони};
        \node[appblock, left=4.2cm of web] (rest) {REST представяне};
        \node[appaccent, right=4.2cm of service] (config) {Конфигурация\\Bean-и и i18n};

        \draw[appline] (client) -- (web);
        \draw[appline] (web) -- (service);
        \draw[appline] (service) -- (repo);
        \draw[appline] (repo) -- (db);
        \draw[appline] (web) -- (view);
        \draw[appline] (view) |- (client);
        \draw[appline] (web) -- (rest);
        \draw[appline] (rest) |- (client);
        \draw[appline] (config) -- (service);
        \draw[appline] (config) -- (web);
    \end{tikzpicture}
    \caption{Високо ниво на архитектурата на приложението}
    \label{fig:high-level-architecture}
\end{figure}

Фигура \ref{fig:high-level-architecture} показва, че приложението следва централен уеб поток, но може да предоставя резултати както чрез HTML визуализация, така и чрез JSON представяне. Това е важен белег за архитектурна гъвкавост, въпреки че системата е доминиращо сървърно рендирана.

\subsection{Пакетна организация на сорс кода}

Основният Java код е разположен в директорията \texttt{src/main/java/com/example/diplomenproekt}. Вътре в нея проектът е разделен на няколко логически под-пакета. Тази организация не е просто въпрос на стил; тя изразява пряко архитектурното разделение на отговорностите.

\begin{table}[H]
    \centering
    \caption{Пакети в Java слоя и тяхното предназначение}
    \label{tab:package-organization}
    \begin{tabularx}{\textwidth}{L{3.5cm}Y}
        \toprule
        Пакет & Основно предназначение \\
        \midrule
        \texttt{config} & Съдържа конфигурационни класове за bean дефиниции и локализация \\
        \texttt{models.entities} & Описва JPA entity обектите на домейна \\
        \texttt{models.bindingModel} & Съдържа модели за приемане и валидация на входни форми \\
        \texttt{models.serviceModel} & Описва междинни модели за пренос на данни към service слоя \\
        \texttt{models.viewModel} & Съдържа модели, предназначени за визуализация в шаблоните или API \\
        \texttt{repositories} & Дефинира интерфейсите за достъп до базата данни чрез Spring Data JPA \\
        \texttt{services} & Съдържа service интерфейси, описващи бизнес възможностите \\
        \texttt{services.impl} & Реализира бизнес логиката и координацията между компонентите \\
        \texttt{web} & Съдържа MVC и REST контролери за входящите HTTP заявки \\
        \bottomrule
    \end{tabularx}
\end{table}

От таблица \ref{tab:package-organization} се вижда, че дори при сравнително малък проект е постигната смислена пакетна сегрегация. Това е едно от силните качества на разработката, защото улеснява навигацията, поддръжката и аналитичното описание на системата.

\subsection{Входна точка и стартиране на приложението}

Входната точка на системата е класът \texttt{DiplomenProektApplication}, маркиран с анотация \texttt{@SpringBootApplication}. Тази анотация обединява механизми за component scanning, auto-configuration и bootstrapping на приложението \cite{springboot}. На практика това означава, че при стартиране Spring открива контролерите, услугите, repository интерфейсите и конфигурационните класове, създава необходимите bean инстанции и подготвя средата за обслужване на заявки.

От архитектурна гледна точка е важно, че входният клас не съдържа бизнес логика. Той изпълнява единствено ролята на стартова точка. Това е добра практика, защото предотвратява смесването на инфраструктурна и домейн логика още на най-високо ниво на приложението.

\subsection{Конфигурационен слой}

В пакета \texttt{config} са разположени два съществени класа: \texttt{ApplicationBeanConfiguration} и \texttt{I18nConfig}. Първият дефинира bean компоненти за \texttt{ModelMapper} и \texttt{PasswordEncoder}. Тук архитектурната идея е ясна: компоненти с инфраструктурна роля, използвани на множество места в системата, трябва да се управляват централизирано от Spring контейнера, а не да се създават ръчно във всеки клас.

\texttt{I18nConfig} от своя страна настройва \texttt{LocaleResolver} и \texttt{LocaleChangeInterceptor}, с което създава базов механизъм за локализация на интерфейса. Това показва, че архитектурата не е ограничена единствено до CRUD логика, а обхваща и допълнителни крос-срезови характеристики на приложението.

Централизираната конфигурация има и друго предимство: когато даден инфраструктурен избор трябва да бъде променен, това става на едно определено място. Така например подмяна на password encoder или промяна на локализационния механизъм не изисква търсене и редакция в множество несвързани класове.

\subsection{Уеб слой: контролери и маршрути}

Уеб слоят е реализиран в пакета \texttt{web}. В него се намират контролерите \texttt{HomeController}, \texttt{UserController}, \texttt{CompanyController}, \texttt{ProductController} и \texttt{MaxPriceRestController}. Разделянето им по функционални области е логично и съответства на основните домейн сценарии на системата.

\texttt{HomeController} обслужва началната точка на приложението и реализира проста логика за пренасочване към каталога при активна сесия. \texttt{UserController} е най-богатият по отношение на потоци компонент, защото управлява регистрацията, входа, изхода, профила, списъка с потребители и потвърждението на електронна поща. \texttt{CompanyController} обслужва формите за въвеждане на фирма. \texttt{ProductController} управлява продуктовия каталог, добавянето на продукт, детайлния изглед, количката, checkout процеса и историята на поръчките. \texttt{MaxPriceRestController} предоставя отделен JSON интерфейс.

От архитектурна гледна точка силна страна на проекта е, че контролерите до голяма степен запазват ролята си на координатори, а не на място за тежка бизнес логика. Те валидират достъпа до маршрута, проверяват резултата от binding операцията, извикват съответна услуга и определят към коя view да премине изпълнението. Това следва добрата практика контролерът да бъде „тънък“ и ориентиран към уеб потока, а не към самото бизнес правило.

\subsection{Service слой: сърце на бизнес логиката}

Service слоят е организиран чрез интерфейси и имплементации. Това дава две важни предимства. Първо, архитектурно се разграничават договорът и конкретната реализация. Второ, системата става по-подходяща за бъдещо тестване и подмяна на конкретни имплементации.

В настоящия проект могат да бъдат разграничени следните основни service области:

\begin{itemize}[leftmargin=*]
    \item потребителски операции -- регистрация, вход, потвърждение, редакция на профил, изчисляване на обща цена в количката;
    \item управление на фирми -- добавяне и извличане на фирми за конкретен потребител;
    \item управление на продукти -- добавяне, категоризация, търсене по идентификатор, проверка за наличност, добавяне в количка;
    \item количка -- съхранение и извличане на записи за избрани продукти;
    \item история -- финализиране на поръчка, изчистване и извличане на минали покупки.
\end{itemize}

Service слоят е ключов за архитектурата, защото именно там се вземат решения за последователността на операциите. Например при checkout процеса е необходимо не просто да се създаде исторически запис, а да се координира обновяването на количката, историята и продуктите. Този тип действия са типичен пример за логика, която не принадлежи нито на контролера, нито на repository интерфейса, а на самия бизнес слой.

\subsection{Persistence слой и repository компоненти}

Persistence слоят е реализиран чрез Spring Data JPA repository интерфейси: \texttt{UserRepository}, \texttt{CompanyRepository}, \texttt{ProductRepository}, \texttt{ShoppingCartRepository} и \texttt{HistoryRepository}. Тяхната основна функция е да предоставят стандартизиран и сравнително декларативен достъп до данните.

Наличието на методи като \texttt{findByEmail}, \texttt{getAllByUser\_Id}, \texttt{findAllByUser\_Id} и \texttt{deleteAllByUser\_Id} показва типичния за Spring Data подход: част от необходимата query логика се изразява директно чрез имената на методите, без ръчно писане на SQL. Това спестява boilerplate код и прави част от persistence слоя по-компактен.

От архитектурна перспектива repository слоят е правилно ограничен до операции по извличане, запис и изтриване. Той не съдържа самостоятелна бизнес логика, а служи като механизъм, чрез който service слоят реализира домейн правилата. Именно това разделение е един от важните признаци за архитектурна последователност.

\subsection{Модели за трансфер на данни между слоевете}

Интересна архитектурна особеност на проекта е използването на няколко типа модели:

\begin{itemize}[leftmargin=*]
    \item entity модели за persistence слоя;
    \item binding модели за входни форми;
    \item service модели за междинен трансфер;
    \item view модели за визуализация и API представяне.
\end{itemize}

Този подход съответства на добрите практики за ограничаване на директната зависимост между базата данни и потребителския интерфейс. Например не е необходимо HTML шаблоните да получават винаги самите entity обекти; вместо това те могат да работят с view модели, съдържащи само необходимите за визуализация полета. Аналогично, при приемане на данни от форма е по-удачно да се използва binding модел с валидационни анотации, отколкото директно persistence обектът.

Използването на \texttt{ModelMapper} подсказва стремеж към намаляване на ръчните преобразувания между тези представяния. Това е полезно от инженерна гледна точка, но не отменя необходимостта от внимателно моделиране на самите класове и полета.

\subsection{Представящ слой: шаблони и статични ресурси}

Визуалният слой на приложението е разположен в \texttt{src/main/resources/templates} и \texttt{src/main/resources/static}. Шаблоните отразяват основните потребителски екрани: начална страница, регистрация, вход, профил, добавяне на фирма, добавяне на продукт, каталог, детайли, количка, история и страници за грешки. Статичните ресурси включват CSS и JavaScript файлове, както и изображения.

Архитектурното значение на този слой се състои в това, че той остава относително изолиран от persistence логиката. Шаблоните работят с вече подготвени атрибути от модела и не съдържат сложни бизнес решения. Това е важен белег за дисциплинирано разделение на отговорностите.

\subsection{Архитектурен поток на ключови операции}

За да бъде по-ясно как слоевете си взаимодействат, е полезно да се разгледа общ модел на няколко основни операции.

\subsubsection{Регистрация}

При регистрация браузърът изпраща POST заявка към потребителския контролер. Контролерът приема binding модел и го предава към потребителската услуга. Услугата проверява дали вече съществува потребител с този имейл, създава entity обект, хешира паролата, записва новия потребител чрез repository слоя и изпраща потвърждаващ имейл чрез инфраструктурен компонент. След това контролерът връща подходящо пренасочване.

\subsubsection{Добавяне на продукт}

Контролерът за продукти валидира входната форма и предава данните към \texttt{ProductsService}. Услугата намира избраната фирма, създава нов продукт, преобразува категорията от низ към \texttt{CategoryEnum} и записва обекта в базата. По този начин контролерът остава концентриран върху маршрута и валидацията, а service слоят поема домейн решението.

\subsubsection{Checkout}

Checkout сценарият е най-показателен за архитектурната координация. Контролерът приема адрес и обща цена, валидира входа и извиква \texttt{HistoryService}. Тази услуга създава запис в историята, асоциира го с потребителя, нулира или изчиства данните за текущата количка и премахва продукти с нулева наличност. Тук ясно се вижда защо service слоят е необходим: операцията изисква координация между няколко persistence компонента и не може да бъде сведена до една repository заявка.

\subsection{Силни страни на архитектурното решение}

От направения анализ могат да бъдат изведени няколко силни страни на архитектурата:

\begin{itemize}[leftmargin=*]
    \item ясна пакетна организация, която отразява логическите роли на класовете;
    \item разграничение между контролери, услуги и repository интерфейси;
    \item използване на отделни модели за различни етапи на обмен на данни;
    \item централизация на инфраструктурни компоненти чрез конфигурационни bean дефиниции;
    \item наличие както на HTML, така и на примерен REST интерфейс.
\end{itemize}

Тези характеристики правят системата подходяща не само за изпълнение на текущата функционалност, но и за последващ анализ и разширение.

\subsection{Архитектурни компромиси и ограничения}

Наред със силните страни е необходимо да бъдат отчетени и определени компромиси. Част от проверките за достъп са реализирани директно в контролерите чрез условие върху сесията, вместо чрез централизирана security архитектура. Това прави системата по-лесна за разбиране, но по-слаба откъм формална сигурност. Някои операции в service слоя показват силна зависимост от конкретната структура на данните, което е очаквано за малък проект, но при по-голямо приложение би следвало да бъде допълнително абстрахирано.

Допълнително, макар архитектурата да е многослойна, тя все още е тясно обвързана с конкретната реализация на сесиен достъп и локална конфигурация. Това е разумно за учебен сценарий, но следва да бъде адресирано при бъдещо развитие към по-зряла среда.

\subsection{Обобщение}

Архитектурата на разглежданата система представлява добре подредена монолитна Spring Boot реализация с ясно разграничени слоеве и пакетна структура. Използваният модел е напълно адекватен за домейн като онлайн магазин с ограничен, но реалистичен обхват. Той позволява логическа проследимост на данните и операциите, създава добра основа за поддръжка и надграждане и подчертава учебно-инженерната стойност на разработката. Следващата глава преминава към по-детайлно разглеждане на домейн модела и начина, по който обектите на системата се картографират към базата данни и към останалите слоеве.
